\input header
\input ras.tex
\def\caratheodory{Carath\'eodory }
\footline={\hss\tenrm 15-\folio\hss}\pageno=1

\centerline{\bf CARATH\'EODORY THEORY}\bigskip

\centerline{\bf \S15. Prime Ends.}\medskip

Let \[U\] be a simply connected subset of \[\hat\C\] such that the complement
\[\hat\C-U\] is infinite. Then the Riemann Mapping Theorem asserts that
there is a conformal isomorphism
$$	\phi:U \buildrel \approx\over\longrightarrow D\,.	$$
In some cases, \[\phi\] will extend to a homeomorphism from the closure
\[\bar U\] onto the closed disk \[\bar D\]. (Compare Figures 1a and 14a,
together with \S16.7.)
However, this is certainly not true in general, since
the boundary \[\partial U\] may be an extremely complicated object.
As an example, Figure 14b shows a region \[U\] such that one point of
\[\partial U\] corresponds to a Cantor set of distinct points of the circle
\[\partial D\], and Figures 14c, 14d show examples
for which an entire interval
of points of \[\partial U\] corresponds to a single point of the circle.
An effective analysis of the relationship between the compact
set \[\partial U\] and the boundary circle \[\partial D\] was carried out
by \caratheodory in 1913, and will be described here.\smallskip

\midinsert
\hbox to \hsize{
\hfil
\hbox{\RasterBox {fig14a.ps} {480} {480} {400}}
\hbox{\RasterBox {fig14b.ps} {480} {480} {400}}
\hfil
}
\vskip -.5in
\hbox to \hsize{
\hfil
\hbox{\RasterBox {fig14c.ps} {480} {480} {400}}
\hbox{\RasterBox {fig14d.ps} {480} {480} {400}}
\hfil
}

\centerline{Figure 14. The boundaries of
four simply connected regions in \[\C\].}
\vfil
\endinsert

The final construction will be purely topological, but we must first
use analytic methods, as in Appendix B,
to prove several lemmas about the existence of short arcs.
Let \[I=(0,\delta)\] be an open interval of real numbers, and let \[I^2\subset
\C\] be the open square, consisting of all \[z=x+iy\] with \[x,\,y\in I\].
We will always use the spherical metric on \[\hat\C\] when measuring
either arclengths or areas. Recall from \S1(6) that this metric has the form
\[ds=\eta(z)|dz|\] with \[\eta>0\]. (In fact
\[\eta(z)=2/(1+|z|^2)\].) Setting \[z=x+iy\], the corresponding spherical
area element on \[\hat\C\] can be written as \[dA=\eta^2(z)
dxdy\], with total area \[\int\!\!\int dA=4\pi<\infty\].

{\QP{\bf 15.1. Lemma.} {\it If \[f:I^2\to V\] is a conformal isomorphism
from the open square onto an open subset of \[\hat\C\], then for
Lebesgue almost every \[x\in I\] the arc \[f(x\times I)\]
has finite spherical arclength.}\medskip}

{\bf Proof.} (Compare B.1 in the Appendix.)
The length \[L(x)\] of this image \[f(x\times I)\]
can be expressed as \[\int_I\, |\eta f'|dy\]
where the function \[\eta\] is to be evaluated at \[f(x+iy)\],
while the area of \[V\] can be expressed as \[A=\int_I\, a(x)dx\] with
\[a(x)=\int_I\, |\eta f'|^2dy\]. By the Schwarz inequality we have
$$      \bigl(\int_I 1\cdot |\eta f'|\;dy\,\bigr)^2\;\le \;\bigl(
	\int_I 1^2dy\bigr)\cdot\bigl(\int_I |\eta f'|^2 dy\bigr)\,, $$
or in other words \[L(x)^2\le \delta a(x)\]. Since the integral of \[a(x)\] is
finite, it follows that \[a(x)\] is finite almost everywhere, hence
\[L(x)\] is finite almost everywhere.\QED

Here is a more quantitative version of this statement. Again let
\[A\] be the spherical area of \[f(I^2)=V\].

{\QP{\bf 15.2. Lemma.} {\it With \[f:I^2\buildrel \approx\over\to
V\] as above, the length \[L(x)\] of the curve \[f(x\times I)\] satisfies
\[L(x)<2\sqrt A\] for more than half of the points \[x\in I\]. Similarly, the
length of \[f(I\times
y)\] is less than \[2\sqrt A\] for more than half of the points
\[y\in I\].}\medskip}


{\bf Proof.} We must show that the Lebesgue measure of
\[S=\{x\in I:L(x)\ge 2\sqrt{A}\}\] satisfies \[\ell(S)<{1\over 2}
\ell(I)={1\over 2}\delta\]. It is clear that \[4A\ell(S)=(2\sqrt A)^2
\ell(S)\le\int_I L(x)^2dx\].
On the other hand, the
discussion above shows that \[\int_I L(x)^2dx\le\delta A\]. Combining these
two inequalities and dividing by \[4A\], we obtain \[\ell(S)\le\delta/4\],
as required.\QED

Now consider a simply connected open set \[U\subset\hat \C\] and some
choice of Riemann map \[\phi:U\to D\], with inverse \[\psi:D\to U\].

{\QP{\bf 15.3. Theorem.} {\it For almost every point \[e^{i\theta}\] of
the circle \[\partial D\] the radial line \[r\mapsto re^{i\theta}\]
maps under \[\psi\] to a curve of finite spherical
length in \[U\]. In particular, the radial limit
$$	\lim_{r\to 1}\; \psi(re^{i\theta})\in\partial U	$$
exists for Lebesgue almost every \[\theta\]. Furthermore, if we fix
\[\theta\], then for Lebesgue almost every \[\theta'\] the radial limit
of \[\psi(re^{i\theta'})\] is different from the radial limit of \[\psi(r
e^{i\theta})\].}\medskip}

\eject
We will say briefly that almost every image curve \[r\mapsto\psi
(re^{i\theta})\]
in \[U\] {\bit lands} at some single point of \[\partial U\], and that
different values of \[\theta\] almost always correspond to distinct
landing points.

{\bf Remark.} The first part of this Lemma is the univalent version of a
basic {\bit Theorem of Fatou}, which
Carath\'eodory used as the starting point of his theory. Fatou's Theorem
says that any bounded holomorphic function on \[D\] has radial limits
in almost all directions, whether or not it is univalent. (See for example
Hoffman, p. 38.) However the univalent case is all that we will need,
and is easier to prove than the general theorem.\smallskip

{\bf Proof of 15.3.}
Map the upper half-plane onto \[D-\{0\}\] by the exponential
map \[w\mapsto e^{iw}\]. In terms of polar coordinates \[r,\,\theta\] in
\[D\], this means that we set \[w=\theta-i\log r\], so that \[e^{iw}=r
e^{i\theta}\]. Applying the argument of 15.1
to the composition \[w\mapsto \psi(e^{iw})\], we see that \[\psi\] has
radial limits in almost all directions \[e^{i\theta}\].

If this map \[\psi\] is bounded, then a
Theorem of F. and M. Riesz asserts that any
given radial limit can occur only for a set of directions \[e^{i\theta}\]
of measure zero. A proof may be found in Appendix A, Theorem A.3. Now for any
univalent \[\psi\], we can reduce to the bounded case in two steps,
as follows. First suppose that
the image \[\psi(D)=U\] omits an entire neighborhood of some point \[z_0\]
of \[\hat\C\]. Then by composing \[\psi\] with a fractional linear
transformation which carries \[z_0\] to \[\infty\], we reduce to the
bounded case. In general, \[\psi(D)\] must omit at least two values,
which we may take to be \[0\] and \[\infty\]. Then
\[\sqrt \psi\] can be defined as a single valued function which omits an
entire open set of points; and we are reduced to the previous case.\QED

Similarly, using the change of coordinates \[z=e^{iw}\] together
with 15.2, we obtain the following. Let \[\psi:D\buildrel\approx\over
\to U\] as above.

{\QP{\bf 15.4. Lemma.} {\it Any point of \[\partial D\] has a nested
sequence of neighborhoods\break
\[{\cal N}_1\supset{\cal N}_2\supset\cdots\] in \[\bar D\] with the following
properties:

\ref $(1)\] each boundary \[\overline{{\cal N}_k}\cap\overline{D-{\cal N}_k}\]
consists of an open arc \[{\cal A}_k\]
in the interior of \[D\] together with two end points in \[\partial D\],

\ref $(2)\] these boundaries \[\bar{\cal A}_k\] are pairwise disjoint,

\ref
$(3)\] both \[{\cal A}_k\] and its image \[\psi({\cal A}_k)\subset U\] have
finite length which tends to zero as \[k\to\infty\], and

\ref $(4)\] each image arc \[\psi({\cal A}_k)\] has two distinct boundary
points in \[\partial U\], and the\break
\hbox{\kern .8em} closures \[\overline{\psi({\cal A}_k)}\]
are pairwise disjoint.}\medskip}

{\bf Proof.} We choose a neighborhood bounded by two short radial line
segments going out to the boundary of \[D\] together with an arc of a
circle \[|z|={\rm constant}\]. Using 15.2 we see that the image of such
a curve \[{\cal A}\] can have length less than any given epsilon, and using
the last part of 15.3 we see that \[{\cal A}\] can be chosen so that the
two endpoints of \[\psi({\cal A})\] in \[\partial U\] are distinct,
and so that successive arcs have distinct endpoints.\QED

Conversely, let us start with an embedded arc \[\alpha:[0,1)\to U\] which
lands at a point \[z_0\] of \[\partial U\]. That is, we suppose
that \[\lim_{t\to 1}\, \alpha(t)=z_0\].
\eject

{\QP{\bf 15.5. Theorem.} {\it Any arc in \[U\] which lands at one point
\[z_0\] of \[\partial U\] corresponds, under the Riemann map,
to an arc in \[D\] which lands at one point of \[\partial D\].
Furthermore, arcs which land at distinct points of \[\partial U\] necessarily
correspond to arcs which land at distinct points of \[\partial D\].}\medskip}

{\bf Proof.}
If the image arc \[t\mapsto \phi(\alpha(t))\] did not land at a single
point of \[D\], then it would have to oscillate back and forth (or round
and round), accumulating
towards some connected set of limit points on \[\partial D\]. It would then
follow, for some set of \[e^{i\theta}\in \partial D\]
of positive measure, that the radial
limit of \[\psi(re^{i\theta})\] could only be \[z_0\], which is impossible
by the Riesz Theorem.

If arcs landing at two distinct points
of \[\partial U\] corresponded to arcs landing at a single point of \[\partial
D\]. Then, using 15.4, we could cut both of these arcs by the image of an arc
which is arbitrarily short and arbitrarily close to the
boundary both in \[D\] and in \[U\]. Evidently this is impossible.\QED

Using these simple facts, \caratheodory showed that we can reconstruct the
topology of the closed disk \[\bar D\] from the topology of the embedding
of \[U\] into \[\bar U\]. There are a number of possible variations on his
basic definitions. (Compare Ahlfors [1973], Epstein, Ohtsuka.)
We will make use of a variation which is fairly close to the original
construction.\smallskip

A {\bit transverse arc} (or ``crosscut'') in \[U\] is a set
\[A\subset\bar U\] which is homeomorphic to the closed interval \[[0,1]\],
and which intersects the boundary \[\partial U\] only in the two end
points of the interval.\smallskip

Note that it is very easy to construct examples of transverse arcs.
For example we can start with any short line segment inside \[U\] and extend
in both directions until it first hits the boundary.

{\QP {\bf 15.6. Lemma.} {\it Any transverse arc \[A\] cuts \[U\] into two
connected components.}\medskip}

{\bf Proof.} The quotient space \[\bar U/\partial U\], in which the boundary
is identified to a point, is evidently homeomorphic to the 2-sphere. Since
\[A\] corresponds to a Jordan curve in this quotient 2-sphere, the
conclusion follows from the Jordan Curve Theorem. (See for example
Munkres.)\QED

It will be convenient to distinguish these two components of \[U-A\]
by choosing some base point \[b_0\in U\].
Then for any transverse arc \[A\] which is disjoint from \[b_0\] we
define the  neighborhood \[N(A)\] which is ``cut off'' by \[A\] to be the
component of \[U-A\] which does not contain \[b_0\].\smallskip

{\bf Definition.}
A {\bit fundamental chain} is a infinite sequence \[A_1\,,\, A_2\,,\,\ldots\]
of disjoint transverse arcs with the property that the corresponding
neighborhoods are nested:
$$	 N(A_1)\supset N(A_2)\supset N(A_3)\supset\cdots\,,	$$
and with the property that the diameter of \[A_i\] tends to zero as
\[i\to\infty\].\smallskip

Note however that the neighborhoods \[N(A_i)\] are not required to become
small as \[i\to\infty\]. (Compare Figure 14c, d.)
By the {\bit impression} (or {\bit support})
of a fundamental chain \[\{A_i\}\] we mean
the intersection of the closures \[\bar N(A_i)\]. It is not difficult to
check that this impression is always a compact connected subset
of \[\partial U\].
Evidently the impression consists of a single point if and only if the
diameter of \[N(A_i)\] tends to zero as \[i\to\infty\].\smallskip

Two such fundamental chains \[\{A_i\}\] and \[\{A'_j\}\] are said to be
{\bit equivalent} if each \[N(A_i)\] contains some
\[N(A'_j)\] and each \[N(A'_j)\] contains some \[N(A_i)\]. An equivalence
class of fundamental chains is also called a {\bit prime end} \[{\cal E}\]
in \[U\], or a point in the {\bit \caratheodory boundary} of \[U\].

Two fundamental chains \[\{A_i\}\] and \[\{A'_j\}\] are said to be
{\bit disjoint} if
\[	N(A_i)\cap N(A'_j)=\emptyset	\]
for some \[i\] and \[j\].

{\QP{\bf 15.7. Lemma.} {\it Any two fundamental chains are either equivalent
or disjoint.}\medskip}

{\bf Proof.} Given \[A_i\] we can choose \[j\] large enough so that the
diameter of \[A'_j\] is less than the distance between \[A_i\] and
\[A_{i+1}\]. It then follows easily that \[N(A'_j)\] must be either
contained in \[N(A_i)\] or disjoint from \[N(A_{i+1})\].\QED

We can now define the {\bit \caratheodory completion} \[\hat U\] of \[U\]
to be the disjoint union of the set \[U\] and the set of all prime ends of
\[U\], topologized as follows. For any transverse arc \[A\subset \bar U-
\{b_0\}\] we define the neighborhood \[{\cal N}_A\] to be
the neighborhood \[N(A)\subset U\] together with the set consisting of all
prime ends \[{\cal E}\] for which some representative fundamental chain
\[\{A_i\}\] satisfies \[A_1=A\]. These neighborhoods \[{\cal N}_A\],
together with the open subsets of \[U\], form a basis for the required
topology.

{\QP{\bf 15.8. Lemma.} {\it In the case of the unit disk \[D\], the
identity map of \[D\] extends uniquely to a homeomorphism between the closure
\[\bar D\subset\C\] and the \caratheodory completion \[\hat D\].}
\medskip}

(Compare 16.7.) In particular, the prime ends of \[D\]
are in one-to-one correspondence with the points of \[\partial D\].
More precisely, the impression of any prime end of \[D\] is a single point of
\[\partial D\], and
each point of \[\partial D\] is the impression of one and only one
prime end. The proof is not difficult.\QED

For an arbitrary simply connected and Hyperbolic
\[U\subset\hat\C\] we now have the following.

{\QP{\bf 15.9. Theorem.} {\it The Riemann map \[\phi:U\to D\] extends
uniquely to a homeomorphism between the completion \[\hat U\] and the
completion \[\hat D\cong\bar D\].}\medskip}

{\bf Proof.} It follows immediately from 15.5 that every transverse arc
in \[U\] corresponds to a transverse arc in \[D\], and that disjoint
transverse arcs correspond to disjoint transverse arcs. Thus every end
in \[U\] gives rise to an end in \[D\]. On the other hand,
given any point of \[\partial D\], we can find a nested sequence of
neighborhoods in \[\bar D\], each bounded by a figure consisting of two
radial segments and an arc which is completely inside \[D\]. As in 15.4,
we can choose these figures so as to correspond to transverse arcs in
\[\bar U\] which are pairwise disjoint, with lengths converging to zero.
The resulting fundamental chain in \[U\] determines the required prime end.
Further details will be left to the reader.\QED
\eject
\footline={\hss\tenrm 16-\folio\hss}\pageno=1


\centerline{\bf \S16. Local Connectivity.}\bigskip

This will be a continuation of the preceeding section, describing
Carath\'eodory's theory in the locally connected case. First a number
of definitions and well known lemmas.
We will always assume that our topological spaces are Hausdorff.
By definition, a topological space \[X\] is {\bit locally connected} at
the point \[x\in X\] if there exist arbitrarily small connected neighborhoods
of \[x\] in \[X\].

Here our ``neighborhoods'' are not required to be open subsets
of \[X\], but only to be sets which contain a (relatively) open set
containing \[x\]. Unfortunately, some authors make a slightly different
definition by considering only open neighborhoods. Let us say
that \[X\] is {\bit openly locally connected} at a point \[x\] if there exist
arbitrarily small connected open neighborhoods of \[x\] in \[X\]. This
is a stronger condition (compare Problem 16-1); however we have the following.

{\QP{\bf 16.1. Lemma.} \it The space \[X\] is locally connected at every point
\[x\in X\] if and only if every open subset of \[X\] is a union of
connected open subsets of \[X\].\medskip}

\noindent If this conditions is satisfied, then evidently
\[X\] is openly locally connected at every point.
Such spaces are simply said to be {\bit locally connected}.
The proof is straightforward, and will be left to the reader.\QED

{\bf 16.2. Remark.} If \[X\] is compact metric, with distance
function \[\rho(x,x')\], then an equivalent condition is that:

{\narrower\smallskip\noindent\it
For any \[\epsilon>0\] there exists \[\delta>0\] so that any two points
with distance\break \[\rho(x,x')<\delta\] are contained in a connected set of
diameter less than \[\epsilon\].\smallskip}

\noindent Again the proof is straightforward, and will be omitted.\smallskip

The space \[X\] is {\bit path connected}$\, $ if there exists a continuous
map from the unit interval \[[0,1]\] into \[X\] which joins any two given
points. It is {\bit arcwise connected}$\,$ if there is a topological embedding
of \[[0,1]\] into \[X\] which joins any two given distinct points.

{\QP{\bf 16.3. Lemma.} {\it A space is path connected if and only if it
is arcwise connected.}\medskip}

\noindent (Recall that all spaces must be Hausdorff.)
Proofs will be deferred until the end of this section.
There is a corresponding concept of {\bit ``locally path connected''} or
equivalently {\bit ``locally arcwise connected''}.
In general, a connected space need not be path connected (Figure 14d). However:

{\QP{\bf 16.4. Lemma.} {\it If a compact metric space \[X\] is locally
connected, then it is locally path connected. Hence every connected
component of \[X\] is actually path connected.}\medskip}

\noindent We will also need the following statement.

{\QP{\bf 16.5. Lemma.} {\it Any continuous image of a compact locally
connected space is compact and locally connected.}\medskip}

\noindent The following is the principal result of this section. Let \[U\subset
\hat\C\] be an open set which is conformally isomorphic to the unit disk.
\eject

{\QP{\bf 16.6. Theorem of Carath\'eodory.} {\it The inverse Riemann map
\[\psi:D\buildrel
\approx\over\to U\] extends continuously to a map from the
closed disk \[\bar D\] onto \[\bar U\] if and only if the boundary \[\partial
U\] is locally connected, or if and only if the complement \[\hat\C-U\]
is locally connected.}\medskip}

\noindent \caratheodory also proved the following.

{\QP{\bf 16.7. Theorem.} {\it If the boundary of \[U\] is a Jordan curve, then
the Riemann map extends to a homeomorphism from the closure \[\bar U\]
onto the closed disk \[\bar D\].}\bigskip}

\noindent The proofs begin as follows.\smallskip

{\bf Proof of 16.3.} Let \[f=f_0:[0,1]\to X\] be any continuous path with\break
\[f(0)\ne f(1)\]. We must construct an embedded arc \[A\subset X\] from
\[f(0)\] to \[f(1)\].
Choose a closed subinterval \[I_1=[a_1\,,\,b_1]\subset [0,1]\] whose
length \[0\le\ell(I_1)=b_1-a_1<1\] is as large as possible, subject to the
condition that \[f(a_1)=f(b_1)\]. Now,
among all subintervals of \[[0,1]\] which are disjoint from \[I_1\],
choose an interval \[I_2=[a_2\,,\,b_2]\]  of maximal length subject to
the condition \[f(a_2)=f(b_2)\].
Continue this process inductively, constructing disjoint subintervals
of maximal lengths \[\;\ell(I_1)\ge\ell(I_2)\ge\cdots\ge 0\;\] subject to
the condition that \[f\]
is constant on the boundary of each \[ I_j\].

Let \[\alpha:[0,1]\to X\] be
the unique map which is constant on each closed interval \[I_j\], and which
coincides with \[f\] outside of these subintervals. Thus \[\alpha(t)\] must
coincide with \[f(t)\] for \[t\in\partial I_j\].
Then it is easy to check
that \[\alpha\] is continuous, and that for each point\break
\[x\in \alpha([0,1])\] the
preimage \[\;\alpha^{-1}(x)\subset[0,1]\;\] is a (possibly degenerate) closed
interval of real numbers. {\it We will show that these conditions,
with \[\alpha\] non-constant,
imply that the image \[A=\alpha([0,1])\subset X\;\] is an embedded arc
joining \[\alpha(0)\] to \[\alpha(1)\].}

Choose a countable dense subset \[\{t_1\,,\, t_2\,,\,\ldots\}\subset[0,1]\].
Define a total ordering of the image \[A\] by specifying
that \[\alpha(s)<\alpha(t)\] if and only if \[\alpha(s)\] and \[\alpha(t)\]
are distinct points with \[s<t\].
An order preserving homeomorphism \[h:[0,1]\buildrel\approx\over
\to A\] can now be constructed inductively as follows.
We must set \[h(0)=\alpha(0)\] and \[h(1)=\alpha(1)\]. For each
dyadic fraction \[0<m/2^k<1\] with \[m\] odd, let us assume inductively that
\[h((m-1)/2^k)\] and \[h((m+1)/2^k)\] have already been defined. Then choose
the smallest index \[j\] so that
$$	h\big({m-1\over 2^k}\big)\;<\;\alpha(t_j)\;<\;h\big({m+1\over2^k}\big)	$$
and set \[h(m/2^k)=\alpha(t_j)\]. It is not difficult to check that the \[h\]
constructed in this way on dyadic fractions extends uniquely to an order
preserving map from \[[0,1]\] to \[A\], which is necessarily a
homeomorphism.\QED

{\bf Proof of 16.4.} Let \[X\] be compact metric and locally connected.
Given \[\epsilon>0\], it follows from 16.2 that we can choose a
sequence of numbers \[\delta_n>0\] so
that any two points with distance \[\rho(x,x')<\delta_n\] are contained in
a connected set of diameter less than \[\epsilon/2^n\]. {\it We will prove
that any two points
\[x(0)\] and \[x(1)\] with distance \[\rho(x(0),x(1))
<\delta_0\] can be joined by a path of diameter at most
\[4\epsilon\].} The plan of attack is as follows. We will choose a sequence
of denominators \[1=k_0<k_1<k_2<\cdots\], each of which divides the next.
Also, for each fraction of the form \[i/k_n\] between 0 and 1 we will
choose an intermediate point \[x(i/k_n)\]. These are to satisfy two
conditions:

{\narrower\smallskip\noindent
(1) Any two consecutive fractions \[i/k_n\] and \[(i+1)/k_n\] must correspond
to points \[x(i/k_n)\] and \[x((i+1)/k_n)\] which have distance less
than \[\delta_n\] from each other.\smallskip}

\noindent
Hence any two such points are contained in a connected set \[C(i,k_n)\]
which has diameter less than \[\epsilon/2^n\]. The second condition is that:

{\QP
(2) Each point of the form \[x(j/k_{n+1})\], where \[j/k_{n+1}\] lies between
\[i/k_n\] and \[(i+1)/k_n\], must belong to this set \[C(i,k_n)\], and
hence have distance less than \[\epsilon/2^n\] from \[x(i/k_n)\]
and from \[x((i+1)/k_n)\].\smallskip}

\noindent The construction of such denominators \[k_n\] and intermediate
points \[x(i/k_n)\], by induction on \[n\], is completely straightforward,
and will be left to the reader. Thus we may assume that \[x(r)\] has been
defined for a dense set of rational numbers \[r\] in the unit interval.

Next we will prove that this densely defined correspondence \[\;r\mapsto x(r)
\;\] is uniformly continuous. Let \[r\] and \[r'\] be any two rational numbers
for which \[x(r)\] and \[x(r')\] are defined. Suppose
that \[|r-r'|<1/k_n\]. Then choosing \[i/k_n\] as close as possible to both
\[r\] and \[r'\],
it is easy to show that \[x(r)\] and \[x(r')\] have distance less
than \[2\epsilon/2^n\] from \[x(i/k_n)\]. Hence they have distance less than
\[4\epsilon/2^n\] from each other. This proves uniform continuity; and it
follows that there is a unique continuous extension \[t\mapsto x(t)\]
which is defined for all \[t\in[0,1]\]. In this way, we have constructed the
required path of diameter \[\le 4\epsilon\] from \[x(0)\] to \[x(1)\].
Thus \[X\] is locally path connected; and the rest of the
argument is straightforward.\QED

{\bf Proof of 16.5.} Let \[f(X)=Y\] where \[X\] is compact and locally
connected. Then \[Y\] is clearly compact, and we must show that it is
locally.connected. Given any point \[y\in Y\] and open neighborhood \[V\subset Y\],
choose a smaller neighborhood \[U\] with \[\bar U\subset V\].
Then \[f^{-1}(V)\] is a union of disjoint connected open sets, and finitely
many of these connected open sets, say \[V_1\,,\,\ldots\,,\,V_q\],
will suffice to cover the compact subset \[K=f^{-1}(\bar U)\]. Let \[N\]
be the union of those \[f(V_i)\] which contain the given point \[y\].
Evidently \[N\] is connected and contained in \[V\]. In order to prove
that \[N\] is a neighborhood of \[y\], let \[L\] be the union of those compact
images \[f(K\cap V_i)\] which do not contain \[y\]. Then \[U-L\] is
an open neighborhood of \[y\] which is contained in \[N\]. Thus \[N\subset V
\] is a connected neighborhood of \[y\], which proves that
\[Y\] is locally connected.\QED

{\bf Proof of Theorem 16.6.} We will always use the spherical metric on
\[\hat\C\]. If either \[\partial U\] or \[\hat\C-U\] is
locally connected, then given \[\epsilon\] we can choose \[\delta\] so
that any two points of distance less than \[\delta\] in \[\partial U\] are
joined by an arc \[A\] of diameter less than \[\epsilon\] in \[\hat\C-U\].
Now suppose that we start with some point \[w_0\in\partial D\]. According to
Lemma 15.4, we can find a neighborhood of diameter less than \[\delta\] in
\[\bar D\] which is bounded by an arc which maps under \[\psi:D\buildrel
\approx\over\to U\] to an arc \[A'\subset U\] which has length less than \[\delta\].
Furthermore, we can choose this arc so that its two endpoints in \[\partial
U\] are distinct. Let \[A\subset\hat\C-U\] be an arc in the complement which
has diameter less than \[\epsilon\], and has the same endpoints. Then the union \[A\cup A'\] is a Jordan
curve of diameter less than \[\epsilon+\delta\] in \[\hat\C\]. Hence it
bounds a region of diameter less than \[\epsilon+\delta\]. Now as \[\epsilon
\to 0\] this region must converge to a point in \[\partial U\], which
we define to be \[\psi(w_0)\]. Thus we have constructed a function from the
closed disk \[\bar D\] to \[\bar U\] which extends \[\psi:D\to U\]. It
is not difficult to check that this extension is continuous.

Conversely, if \[\psi\] extends continuously over \[\bar D\], the the boundary
\[\partial U\] is a continuous image of the circle \[\partial D\]. Hence
it follows from Lemma 16.5 that \[\partial U\] is locally connected. To
see that \[\hat\C-U\] must also be locally connected, let \[N\] be an
arbitrarily small connected neighborhood of a point \[z\] within \[\partial U\].
Choose \[\delta>0\] so that the \[\delta$-neighborhood of \[z\] within
\[\partial U\] is contained in \[N\]. Then the \[\delta$-neighborhood of \[z\]
within \[\hat\C-U\], together with \[N\], is clearly also connected. Therefore
\[\hat\C-U\] is locally connected.\QED

{\bf Proof of Theorem 16.7.} If \[\partial U\] is a Jordan curve, that is
a homeomorphic image of the circle, then we certainly have a continuous
extension \[\psi:\bar D\to \bar U\] by the preceeding theorem. Suppose that
two distinct points of the circle \[\partial D\] mapped to the same point \[z_0\]
of \[\partial U\]. Then the straight
line segment \[A\] joining these two points in \[\bar
D\] would map to a Jordan curve \[\Gamma\subset\hat\C\] under \[\psi\].
Such a Jordan curve must separate \[\hat\C\] into two components, of
which \[\Gamma\] is the common boundary. Note that
the two components of \[D-A\] must map into the two distinct components
of \[\hat\C-\Gamma\] under \[\psi\]. Since
\[\Gamma\] intersects the Jordan curve \[\partial U\] in a single point
\[z_0\], it follows that the entire connected set \[\partial U-\{z_0\}\]
must be contained in just one of the two connected components of
\[\hat\C-\Gamma\].
But this means that one of the two components of \[D-A\] must
map into the other component of \[\hat\C-\Gamma\], which is disjoint from
\[\partial U\]. Hence all of the corresponding  boundary points in
\[\partial D\] can only map to the single point
\[z_0\in\Gamma\cap\partial U\] under \[\psi\]. This is impossible by
the Riesz Theorem A.3 of Appendix A.\QED\lln

{\bf Problem 16-1.} Let \[X\subset\C\] be the compact connected set which
is obtained from the unit interval \[[0,1]\] by drawing line segments from
\[1\] to the points \[{1\over 2}(1+i/n)\] for\break
\[n=1\,,\,2\,,\,3\,,\,\ldots\]
and then adjoining the successive images of this configuration under the map
\[z\mapsto z/2\]. (Figure 15.)
Show that \[X\] is locally connected at the origin,
but not openly locally connected.\vfil
\midinsert
\insertRaster wbroom.ps pixels 1152 by 450 scaled 300
\bigskip
\centerline{Figure 15. The witch's broom.}
\endinsert
\vfil

% \rightline{Milnor, Conformal Dynamics Lectures \S15,16 draft, 1-21-90}
\end





